% ============================================================
% preambulo-estetica.tex
% ------------------------------------------------------------
% CAPA REEMPLAZABLE. ES LA QUE UNA EDITORIAL CAMBIA POR LA SUYA
% CUANDO QUIERE MANTENER SU IDENTIDAD VISUAL. NO DEBE DEFINIR
% NINGUNO DE LOS ENTORNOS DEL CONTRATO.
% ------------------------------------------------------------
% ESTE ARCHIVO ES EL ESQUELETO. LOS BLOQUES DE DISEÑO QUE YA
% EXISTEN EN EL preambulo.tex EN USO —titlesec, titletoc,
% estilos de página, diseño de notas al pie, newfloat,
% imakeidx, glossaries, froufrou, siunitx— SE PEGAN EN EL
% APARTADO 8. NO ESTÁN ACÁ PORQUE SON DECISIONES DE DISEÑO
% PROPIAS Y NO CORRESPONDE REESCRIBIRLAS.
% ============================================================
% LAS CUATRO PRESTACIONES
% ------------------------------------------------------------
% El derivado emite PRESTACIONES y no un «modo». Así, agregar
% un estado editorial nuevo no obliga a tocar esta capa, y cada
% bandera dice qué hace en vez de cómo se llama.
%
%   estado_libro     showframe  homeoarchy  enlaces  cruces
%   borrador             si         si         si      no
%   en_correccion        no         si         si      no
%   en_produccion        no         no         si      no
%   publicado            no         no         no      si
%
% LAS CUATRO TOCAN OPCIONES DE PAQUETE, QUE NO SE PUEDEN
% SUPERPONER DESPUES DE CARGADO EL PAQUETE. POR ESO SON
% BANDERAS DEL DERIVADO Y NO REEMPLAZOS DE TEXTO.
% LA EXCEPCION ES hyperref: SE CARGA UNA SOLA VEZ EN EL
% CONTRATO Y LOS COLORES VAN POR \hypersetup, QUE SI SE
% SUPERPONE.
% ============================================================
% CORRECCIONES APLICADAS RESPECTO DEL PREAMBULO EN USO
%   - FUERA \usepackage[utf8]{luainputenc}: LuaLaTeX CONSUME
%     UTF-8 NATIVAMENTE (SC-02).
%   - FUERA \usepackage{svg}: EXIGE Inkscape Y --shell-escape,
%     Y EN LIBROS SOLO SE USAN PNG, JPG Y PDF.
%   - \renewcommand{\footnoterule} ESTABA DEFINIDO DOS VECES
%     CON EL MISMO CUERPO. QUEDA UNO.
%   - babel LO CARGA EL CONTRATO, ANTES DE csquotes Y biblatex,
%     QUE ES EL ORDEN QUE PIDE LA DOCUMENTACION.
%   - luatexja SOLO SI EL LIBRO LO NECESITA.
% ============================================================

\usepackage{ifthen}
\usepackage{calc}

% ============================================================
% 1. GEOMETRIA
% ------------------------------------------------------------
% Los valores vienen del derivado.
%
% showframe ES OPCION DE PAQUETE Y NO SE PUEDE SUPERPONER: SE
% PASA CON \PassOptionsToPackage PARA NO DUPLICAR LA LISTA
% ENTERA DE OPCIONES EN DOS RAMAS DE UN CONDICIONAL.
%
% Hoy los cuatro margenes son iguales en los 45 formatos de la
% base, asi que inner/outer da lo mismo que centering; se usan
% igual para que el dia que se carguen asimetricos no haya que
% tocar nada. En un libro encuadernado el margen interior suele
% necesitar mas aire que el exterior.
% ============================================================
\ifgbshowframe
  \PassOptionsToPackage{showframe}{geometry}
\fi

\usepackage[
  paperwidth=\gbanchopagina,
  paperheight=\gbaltopagina,
  inner=\gbmargeninterior,
  outer=\gbmargenexterior,
  top=\gbmargensuperior,
  bottom=\gbmargeninferior,
  includehead,
  includefoot,
  headsep=14pt,
  footskip=24pt
]{geometry}

% ------------------------------------------------------------
% CRUCES DE CORTE
% ------------------------------------------------------------
% La hoja de montaje es mayor que el formato final: el derivado
% la calcula sumando el margen de cruces por lado. Las cruces
% solo necesitan lugar, y ese lugar es el mismo para cualquier
% formato.
%
% crop DECLARA LA HOJA DE MONTAJE Y geometry EL FORMATO FINAL:
% SON DOS PARES DE MEDIDAS INDEPENDIENTES, NO UNA DERIVADA DE
% LA OTRA. VERIFICADO: crop EXPANDE MACROS EN SUS OPCIONES
% keyval, ASI QUE NO HACE FALTA \crop[...] APARTE.
% ------------------------------------------------------------
\ifgbcruces
  \usepackage[width=\gbanchomontaje,height=\gbaltomontaje,cam,center]{crop}
\fi

% ============================================================
% 2. IDIOMA
% ------------------------------------------------------------
% babel lo carga el contrato, antes de csquotes y biblatex.
% Aca solo se cambian las cadenas, que si son diseno.
% ============================================================
\addto\captionsspanish{\renewcommand{\contentsname}{Sumario}}

% ============================================================
% 3. TIPOGRAFIA
% ------------------------------------------------------------
% Las fuentes son de la estetica, no de la base. Por eso cuando
% una editorial reemplaza esta capa, cambia las fuentes y el
% colofon las nombra correctamente sin que haya que tocar nada
% mas: el contrato solo las lee.
% ============================================================
\usepackage{fontspec}
\usepackage[final,babel=true]{microtype}

\def\gbfuentetexto{Libertinus Serif}
\def\gbfuentetitulos{IBM Plex Sans Condensed}
\def\gbfuentemono{IBM Plex Mono}

\setmainfont{\gbfuentetexto}[
  Numbers={OldStyle,Proportional},
  Ligatures={TeX,Common,Discretionary},Scale=1.18]
\setsansfont{\gbfuentetitulos}[
Scale=MatchLowercase,
Ligatures=TeX,
Extension=.otf,
UprightFont=*-Regular,
ItalicFont=*-Italic,
BoldFont=*-SemiBold,
BoldItalicFont=*-SemiBoldItalic]
\setmonofont{\gbfuentemono}[
Scale=0.91,
Extension=.otf,
UprightFont=*-Regular,
ItalicFont = IBMPlexMono-Italic.otf,
BoldFont = IBMPlexMono-Bold.otf,
BoldItalicFont = IBMPlexMono-BoldItalic.otf]

%\setsansfont{\gbfuentetitulos}[Scale=MatchLowercase,Ligatures=TeX]
%\setmonofont{\gbfuentemono}[Scale=0.91]

\usepackage{unicode-math}
\setmathfont{Libertinus Math}[Scale=MatchLowercase]

% SOPORTE CJK SOLO SI EL LIBRO LO NECESITA: luatexja PESA EN
% CADA COMPILACION Y CASI NINGUN LIBRO LO USA.
\ifgbcjk
  \usepackage{luatexja-fontspec}
  \setmainjfont{FandolSong}
\fi

\renewcommand{\normalsize}{\fontsize{10pt}{14pt}\selectfont}
\normalsize

% ============================================================
% 4. LISTAS, NOTAS Y FILETES
% ============================================================
% REEMPLAZA EL APARTADO 4 DE preambulo-estetica.tex
%
% QUÉ CAMBIÓ RESPECTO DEL ESQUELETO, Y POR QUÉ
% ------------------------------------------------------------
% 1. SE SACÓ hang DE footmisc
%
%    hang SANGRA EL TEXTO DE LA NOTA HACIA ADENTRO PARA HACERLE
%    LUGAR A LA MARCA. LO QUE SE BUSCA ES LO CONTRARIO: LA MARCA
%    COLGADA EN EL MARGEN Y EL TEXTO OCUPANDO EL ANCHO COMPLETO
%    DE LA CAJA. ESO LO PRODUCE \deffootnote DE scrextend.
%
%    hang Y \deffootnote GOBIERNAN LA MISMA SANGRÍA POR
%    MECANISMOS DISTINTOS, ASÍ QUE CARGAR LOS DOS DEJA EL
%    RESULTADO A MERCED DE CUÁL SE APLIQUE ÚLTIMO. QUEDA UNO.
%
%    bottom Y stable SÍ SE CONSERVAN: FIJAN LAS NOTAS AL PIE DE
%    LA CAJA Y LAS PERMITEN EN LOS TÍTULOS, Y NINGUNO TOCA LA
%    SANGRÍA.
%
% 2. SE SACARON LOS DOS \patchcmd SOBRE \@footnotetext
%
%    VERIFICADO: LOS DOS FALLAN. scrextend REDEFINE
%    \@footnotetext, ASÍ QUE CUANDO EL PARCHE BUSCA \footnotesize
%    O \@MM ESOS TOKENS YA NO ESTÁN EN LA DEFINICIÓN.
%
%    EL PRIMERO SE REEMPLAZA POR \setkomafont{footnote}, QUE ES
%    EL MECANISMO PROPIO DE KOMA Y NO DEPENDE DE PARCHEAR NADA.
%    VERIFICADO: CON \small LA MISMA NOTA OCUPA DOS PÁGINAS Y CON
%    \footnotesize UNA.
%
%    EL SEGUNDO SE REEMPLAZA POR \interfootnotelinepenalty, EL
%    PARÁMETRO DEL KERNEL QUE GOBIERNA ESE CORTE. ADVERTENCIA
%    HONESTA: NO PUDE DEMOSTRAR QUE CAMBIE NADA. UNA NOTA DE
%    CATORCE PÁRRAFOS SE PARTIÓ IGUAL CON EL PARÁMETRO Y SIN ÉL.
%    QUEDA PORQUE ES EL MECANISMO DOCUMENTADO Y NO HACE DAÑO,
%    PERO SI EL COMPORTAMIENTO CON NOTAS LARGAS NO ES EL DEL
%    PREÁMBULO LEGACY, ESTE ES EL LUGAR A REVISAR.
% ============================================================

\usepackage{enumitem}
\setlist{nosep,topsep=4pt}

% ------------------------------------------------------------
% NOTAS AL PIE
% ------------------------------------------------------------
\usepackage[bottom,stable]{footmisc}
\usepackage{scrextend}

% SEPARACIÓN ENTRE NOTAS
\setlength{\footnotesep}{10pt}

% FILETE SEPARADOR: 0,5 pt DE GROSOR Y 40% DEL ANCHO DE COLUMNA
\renewcommand{\footnoterule}{%
	\kern-3pt
	\hrule height 0.5pt width 0.4\columnwidth
	\kern6pt}

% NUMERACIÓN: ENTRE CORCHETES, PALO SECO, CUERPO scriptsize
\renewcommand*{\thefootnote}{\scriptsize\sffamily{[\arabic{footnote}]}}

% CUERPO DE LA NOTA
\setkomafont{footnote}{\small}

% ------------------------------------------------------------
% MARCA COLGADA EN EL MARGEN
% ------------------------------------------------------------
% \llap SACA LA MARCA FUERA DE LA CAJA HACIA LA IZQUIERDA, ASÍ EL
% TEXTO DE LA NOTA EMPIEZA EN EL MARGEN Y OCUPA EL ANCHO COMPLETO.
%
% LOS TRES ARGUMENTOS DE \deffootnote SON: SANGRÍA DE LA PRIMERA
% LÍNEA, SANGRÍA DE LAS SIGUIENTES, Y CÓMO SE COMPONE LA MARCA.
% ------------------------------------------------------------
\newcommand*\gbmarcanotaespacio{0em}

\deffootnote{\gbmarcanotaespacio}
{\parindent}
{%
	\makebox[\gbmarcanotaespacio][r]{\llap{\thefootnotemark\quad}}%
}

% ------------------------------------------------------------
% NOTAS LARGAS PARTIDAS ENTRE PÁGINAS
% ------------------------------------------------------------
% Penalización de partir una nota entre dos páginas. Por omisión
% LaTeX pone 100 en la clase book, y el valor alto que usa el
% kernel en otros contextos hace que una nota muy larga desplace
% texto en lugar de partirse.
%
% VER LA ADVERTENCIA DEL ENCABEZADO: NO PUDE MEDIR DIFERENCIA.
% ------------------------------------------------------------
\interfootnotelinepenalty=100

% ------------------------------------------------------------
% RESTO DEL APARTADO
% ------------------------------------------------------------
\usepackage[labelfont=bf,font=small,labelsep=period,format=plain]{caption}
\usepackage{needspace}

\setcounter{tocdepth}{1}
\setcounter{secnumdepth}{3}
\renewcommand{\thepart}{\arabic{part}}

\raggedbottom
\clubpenalty=10000
\widowpenalty=10000

% ============================================================
% 5. MARCAS TIPOGRAFICAS DE TRABAJO
% ------------------------------------------------------------
% impnattypo marca homeoarquias y problemas de principio y fin
% de linea. La opcion draft es opcion de paquete y no se puede
% superponer: por eso la carga va condicionada y no comentada
% a mano.
% ============================================================
\ifgbhomeoarchy
  \IfFileExists{impnattypo.sty}{%
    \usepackage[hyphenation,homeoarchy,homeoarchywordcolor=orange,
                homeoarchycharcolor=orange,draft]{impnattypo}}{%
    \PackageWarning{gbpublisher}{impnattypo no instalado:
      se pierden las marcas tipograficas de trabajo}}
  \overfullrule=5pt
\fi

% ============================================================
% 6. COLOR DE LOS ENLACES
% ------------------------------------------------------------
% Solo el color. hyperref lo carga preambulo-hyperref.tex, que
% va después de esta capa y lee \ifgbenlaces para decidir entre
% enlaces coloreados y hidelinks.
% ============================================================
\usepackage{xcolor}
\definecolor{gbenlace}{RGB}{160,0,120}

% ============================================================
% 7. CONSTANCIA EN EL LOG
% ------------------------------------------------------------
% Deja escrito en el .log con que prestaciones se compuso.
% Cuando un PDF sale distinto de lo esperado, esta linea dice
% por que sin tener que abrir el derivado.
% ============================================================
\typeout{[gbpublisher] prestaciones:
  showframe=\ifgbshowframe si\else no\fi,
  homeoarchy=\ifgbhomeoarchy si\else no\fi,
  enlaces=\ifgbenlaces si\else no\fi,
  cruces=\ifgbcruces si\else no\fi}

% ============================================================
% 8. AQUI VAN LOS BLOQUES DE DISENO DEL PREAMBULO EN USO
% ------------------------------------------------------------
%   - titlesec y los estilos de titulo
%   - titletoc: partes, capitulos y secciones del sumario
%   - titleps: estilos de pagina y folios corridos
%     (\gbtituloabreviado ES LO QUE CORRESPONDE EN EL FOLIO
%      CORRIDO, Y HOY EL PREAMBULO NO LO CONSUME)
%   - newfloat: el entorno flotante "imagen"
%   - imakeidx + xindy: names (lo llena esindex desde biblatex),
%     onomastico y concepto (los carga el editor)
%   - glossaries: siglas y glosario
%   - siunitx, froufrou, qrcode, pdfpages, rotating
% ============================================================

% Cita con cambio de márgenes y tamaño tipografico
\renewenvironment{quote}
{\normalsize\list{}{\sf\leftmargin=14pt \rightmargin=0pt}%
	\item\relax}
{\endlist}

% Configuración de los índices
\usepackage[xindy]{imakeidx}
\makeindex[name=names, title={}, intoc=false]
\makeindex[name=onomastico, title={}, intoc=false]
\makeindex[name=concepto, title={}, intoc=false]


% ============================================================
% PERSONALIZACIÓN DE LOS ÍNDICES DE PARTES, CAPÍTULOS Y SECCIONES
% ------------------------------------------------------------
% Paquete requerido:
\usepackage{titletoc}
%
% OBJETIVO GENERAL:
% - Controlar la apariencia de los índices (Tabla de Contenidos)
%   para partes, capítulos y secciones.
% - Ajustar sangrías, tipografía, alineación y estilo de los
%   números de página.
% ============================================================

% ------------------------------------------------------------
% PARTES
% ------------------------------------------------------------
% \titlecontents{part}[<left>]{<above-code>}{<numbered-entry-format>}
%                        {<numberless-entry-format>}{<filler-page>}
%
% Argumentos usados:
%   [0em]                    -> Sangría izquierda de la entrada.
%   {\addvspace{5pt}\sf\bfseries\normalsize\selectfont\filright}
%                             -> Espacio vertical antes de la entrada (5pt),
%                                fuente sans serif, negrita, tamaño normal,
%                                alineación a la derecha (\filright).
%   {\contentslabel[\thecontentslabel]{2.5pc}}
%                             -> Formato de la etiqueta numérica (ej. Parte I)
%                                con ancho de 2.5pc.
%   {}                        -> Formato de entrada sin número.
%   {}                        -> No hay reglas ni página aquí.
\titlecontents{part}[0em]
{\addvspace{5pt}\sf\bfseries\normalsize\selectfont\filright}
{\contentslabel[\thecontentslabel]{2.5pc}}
{}
{}

% ------------------------------------------------------------
% CAPÍTULOS
% ------------------------------------------------------------
% Formato general:
%   - Sangría: 1.5pc
%   - Espacio vertical: 0.4em antes de cada capítulo
%   - Fuente: sans serif, alineación a la derecha (\filright)
%   - Número del capítulo alineado con margen de 1.5pc
%   - Separador de puntos hasta el número de página (\titlerule)
%
% Argumentos:
%   [1.5pc] -> sangría izquierda del capítulo
%   {\addvspace{.4em}\sf\selectfont\filright} -> formato general de la entrada
%   {\contentslabel{1.5pc}} -> ancho reservado para número
%   {\hspace*{-1.5pc}} -> corrección si no tiene número
%   {\titlerule*[1pc]{.}\contentspage[\hspace*{-4pc} {\rm\small\thecontentspage}]}
%      -> regla de puntos de relleno hasta el número de página,
%         página en tamaño pequeño y ligeramente desplazada
\titlecontents{chapter}[1.5pc]
{\addvspace{.4em}\sf\selectfont\filright}
{\contentslabel{1.5pc}}
{\hspace*{-1.5pc}}%
{\titlerule*[1pc]{.}\contentspage[\hspace*{-4pc} {\rm\small\thecontentspage}]}%
[]

% ------------------------------------------------------------
% SECCIONES
% ------------------------------------------------------------
% Formato general:
%   - Sangría: 4.5pc
%   - Fuente: \small, alineación a la derecha
%   - Número de sección: 2.5pc
%   - Ajuste de sangría si no hay número: -2.5pc
%   - Separador de puntos hasta el número de página (\titlerule)
\titlecontents{section}[4.5pc]
{\small\filright}                  % Formato general
{\contentslabel{2.5pc}}            % Número con ancho 2.5pc
{\hspace*{-2.5pc}}                 % Si no tiene número
{\titlerule*[1pc]{.}\contentspage} % Regla de puntos hasta la página


% ============================================================
% CONFIGURACIÓN DE ESTILOS DE TÍTULOS CON titlesec
% ------------------------------------------------------------
% Paquete requerido:
\usepackage[sf,bf,compact,topmarks,calcwidth,pagestyles,clearempty,newlinetospace]{titlesec}
%
% OBJETIVO GENERAL:
% - Personalizar la apariencia de partes, capítulos, secciones,
%   subsecciones, subsubsecciones, párrafos y subpárrafos.
% - Controlar fuente, tamaño, alineación, sangrías y espaciado.
% ============================================================

% ============================================================
% DISEÑO DE PARTES
% ------------------------------------------------------------
% Redefinición de \@part y \@spart para controlar:
% - Numeración de partes (si corresponde)
% - Entrada en el TOC
% - Centrado del título y subtítulo de la parte
% - Tamaño de fuente y estilo tipográfico
%
% \@part[#1]#2:
%   - #1: título corto para TOC y encabezado
%   - #2: título completo para mostrar en la página
%   - Centrado, fuente sans serif (\sf), \LARGE para nombre de parte
%     y \Large para título completo
%
% \@spart#1:
%   - Para partes no numeradas
%   - Solo muestra el título centrado, \LARGE, sans serif
%
\makeatletter
\def\@part[#1]#2{%
	\ifnum \c@secnumdepth >-2\relax
	\refstepcounter{part}%
	\addcontentsline{toc}{part}{Parte \thepart\hspace{1em}#1}%
	\else
	\addcontentsline{toc}{part}{#1}%
	\fi
	\markboth{}{}%
	{\centering
		\interlinepenalty \@M
		\ifnum \c@secnumdepth >-2\relax
		\sf\LARGE\selectfont \partname~\thepart
		\par\vskip 20\p@%
		\fi
		\sf\Large\selectfont
		#2\par}%
	\@endpart
}

\def\@spart#1{%
	\markboth{}{}%
	{\centering
		\interlinepenalty \@M
		\sf\LARGE\selectfont
		#1\par}%
	\@endpart
}
\makeatother

% ============================================================
% DISEÑO DE CAPÍTULO
% ------------------------------------------------------------
% \titleformat{\chapter}[display]{...}{...}{sep}{before-code}[after-code]
% - [display]: título en bloque separado (no inline)
% - \Large: tamaño base del título
% - \vspace{-2.5cm}: ajusta altura vertical (sube el título)
% - \centering: centrado
% - \textsc{\MakeLowercase\chaptertitlename}: 'Capítulo' en versalitas
% - \thechapter: número del capítulo
% - 1.5cm: separación entre número y texto
% - \filright\sf\LARGE: alineación a la derecha, fuente sans serif, tamaño LARGER
\titleformat{\chapter}[display]
{\Large}
{\centering{\textsc{\MakeLowercase\chaptertitlename}}~\thechapter}
{1.5cm}
{\filright\sf\LARGE}
[]
% ------------------------------------------------------------
% ESPACIADO DEL TÍTULO DE CAPÍTULO
% ------------------------------------------------------------
% \titlespacing*{\chapter}{izquierda}{antes}{después}
%
% ANTES = 0pt HACE QUE «capítulo N» ARRANQUE EN EL BORDE SUPERIOR
% DE LA CAJA. LA CLASE book RESERVA 50pt POR OMISIÓN.
%
% LA VERSIÓN CON ASTERISCO SUPRIME LA SANGRÍA DEL PRIMER PÁRRAFO
% DESPUÉS DEL TÍTULO.
%
% NO SE USA \vspace NEGATIVO DENTRO DEL \titleformat: ESO EMPUJA
% HACIA ARRIBA DESPUÉS DE QUE LaTeX YA RESERVÓ EL ESPACIO, Y EL
% RESULTADO DEPENDE DEL CONTEXTO. titlesec TIENE EL PARÁMETRO
% PARA DECIRLO DIRECTO.
% ------------------------------------------------------------
\titlespacing*{\chapter}{0pt}{0pt}{30pt}

% ============================================================
% DISEÑO DE SECCIÓN
% ------------------------------------------------------------
% Fuente: sans serif, negrita, tamaño grande (\large), alineación a la derecha
\titleformat{\section}[hang]
{\sf\bfseries\large\raggedright}
{\thesection}{.5em}{}[]
\titlespacing{\section}
{\parindent}{18pt plus 0.5pt minus 0.5pt}{6.75pt}

% ============================================================
% DISEÑO DE SUBSECCIÓN
% ------------------------------------------------------------
\titleformat{\subsection}[hang]
{\sf\large\raggedright}
{\thesubsection}{.5em}{}[]
\titlespacing{\subsection}
{\parindent}{18pt plus 0.5pt minus 0.5pt}{6.75pt}

% ============================================================
% DISEÑO DE SUBSUBSECCIÓN
% ------------------------------------------------------------
\titleformat{\subsubsection}[hang]
{\rm\bfseries\normalsize\raggedright}
{\thesubsubsection}{.5em}{}[]
\titlespacing{\subsubsection}
{\parindent}{18pt plus 0.5pt minus 0.5pt}{6.75pt}

% ============================================================
% DISEÑO DE PÁRRAFO
% ------------------------------------------------------------
% \paragraph: run-in (alineado inline, no bloque)
% - Negrita
% - Sin número
% - Raya al final del título (---)
\titlespacing{\paragraph}
{0pt}                % margen izquierda
{6.75pt plus 0.5pt minus 0.5pt} % espacio antes
{0pt}                % espacio después
\titleformat{\paragraph}[runin]
{\bfseries}{}{0em}{}[\mbox{ --- }]

% ============================================================
% DISEÑO DE SUBPÁRRAFO
% ------------------------------------------------------------
% \subparagraph: estilo colapsado, centrado
% - Fuente sans serif, negrita, tamaño \large
\titlespacing{\subparagraph}
{\parindent}{18pt plus 0.5pt minus 0.5pt}{6.75pt}
\titleformat{\subparagraph}[hang]
{\sf\bfseries\large\centering}
{\thesubparagraph}{.5em}{}[]

